% Sections 10.2-10.7, from lectures/L10/appendix/a_can_frame_format.md.

\section{Introduction to CAN}
\label{c10:sec:can}

CAN was designed for automotive and industrial systems: many independent controllers (engine,
brakes, instrument cluster, sensors) exchanging short messages reliably, with no single point of
failure.

Point-to-point wiring stops scaling almost immediately. The framed protocol in
\secref{c10:sec:framing} assumed a single shared link; one controller talking to an instrument
cluster, an ABS unit and three sensors would need five separate links, each with its own cabling,
connectors and transceiver. A single shared bus replaces all of them: every node connects to the
same pair of wires, and every node can both listen and transmit on it.

\subsection*{The physical bus}
A real CAN bus is a twisted differential pair, \code{CAN_H} and \code{CAN_L}. A node signals a bit by
driving the two wires apart (\textbf{dominant}) or letting them float together through a resistor
(\textbf{recessive}). An external transceiver chip handles the analogue layer.

This book works entirely above it: one logic value per bit, \code{0} for dominant and \code{1} for
recessive. That is exactly what a transceiver presents to, and accepts from, the digital logic on
either side.

\subsection*{Dominant and recessive}
The whole arbitration scheme rests on a single property:
\begin{itemize}
  \item \textbf{Recessive (\code{1})} is the bus's idle state. Every node releases the bus and the
    transceivers' bias networks pull both wires to the same level.
  \item \textbf{Dominant (\code{0})} is asserted by actively driving the wires apart. If
    \emph{any} node drives dominant while the rest are recessive, the bus reads dominant, regardless
    of how many nodes are involved or which asserted it first.
\end{itemize}

That is exactly a \textbf{wired-AND}: treat every node's output as an input to a bus-wide AND gate,
and \code{0} dominates. Real transceivers implement it in hardware with open drain; the system
testbench in \chapref{c17:ch} models identical behaviour digitally by combining every node's pair
\code{(tx_bus, bus_en)} in the same way.

\subsection*{Multi-master, and why that requires arbitration}
Most embedded buses answer the question ``who may transmit?'' by decree: SPI has a single master,
and I2C, which does define a wired-AND arbitration much like the one below, is also nearly always
run with a single master. CAN makes multi-master the normal case: \textbf{any node may transmit as
soon as the bus is idle}, and if two start at the same instant the bus itself decides who carries
on. There is no retransmission back-off and a collision is not a failure at all. This is CAN's
signature idea, \textbf{arbitration}.

The mechanism in brief; the worked example comes in \secref{c10:sec:arbitration}:
\begin{enumerate}
  \item Every message starts with its identifier, sent bit by bit, MSB first.
  \item Every transmitting node also monitors the bus while it sends each identifier bit.
  \item As long as the bit it sent matches what the bus reads, it goes on transmitting.
  \item The moment a node transmits recessive (\code{1}) but observes dominant (\code{0}), some
    other node is transmitting a \code{0} where it itself sent a \code{1}. It has lost arbitration:
    it stops transmitting immediately and becomes a receiver. The node that carried on driving
    dominant wins, without having noticed anything.
\end{enumerate}

Two consequences of that:
\begin{itemize}
  \item \textbf{Lower identifiers always win}, since they carry more leading dominant bits. Hence
    ``priorities'' in CAN folklore.
  \item \textbf{Every identifier has to be unique.} CAN puts that requirement on the network
    designer; the controller does not enforce it in hardware.
\end{itemize}

\begin{aside}
\textbf{A simplification to flag early.} Step 4 describes a \emph{production} controller, and
seamlessly becoming a receiver mid-frame is part of why real CAN treats a collision as a
non-event. This book's \code{can_controller} (\chapref{c16:ch}--\ref{c17:ch}) stops driving at
exactly the same moment, which is the part that matters to the bus, but then abandons the frame and
sets \code{error} instead of decoding the winner's message. \Secref{c19:sec:gaps} returns to it in
its list of differences from a production controller.
\end{aside}

\Canref{1} takes the bus on its own: the topology, the dominant and recessive asymmetry, the
wired-AND, the two termination resistors, and the relationship between bit rate and bus length
\textemdash{} 1~Mbit/s reaches about 40 m, and the reason is arbitration, which requires a bit to
travel out to the bus's furthest node and back before the bit is over.

% ------------------------------------------------------------------------------------------
\section{Broadcast, not addressing}
\label{c10:sec:broadcast}

The framed protocol in \secref{c10:sec:framing} needed DST and SRC because more than one node could
listen. CAN faces the same problem on the same kind of bus and answers it entirely differently:
\textbf{a CAN frame carries no destination address at all.}

Every node sees every frame. The identifier winning arbitration names no receiver; it names the
\emph{kind of message}, broadcast to everyone and relevant to whoever cares. Every node decides for
itself whether an identifier means anything, typically by comparing it against a small set it is
interested in, in software or in the controller's filtering hardware. A node that is not interested
simply discards the frame; nothing at the protocol level steered it away.

\paragraph{Example}
Suppose a design reserves identifiers by convention,
$\mathit{StatusReq} = \code{0x300} + \mathit{nodeId}$ and
$\mathit{StatusResp} = \code{0x400} + \mathit{nodeId}$. Sensor node 3 is asked for status on
\code{0x303} and replies on \code{0x403}. Every node receives the frame \code{0x303}; only the one
recognising it as its own request identifier replies. No field said ``this is for node 3''.

It is a genuinely different design choice. The protocol in \secref{c10:sec:framing} scales addressing
by widening DST and SRC as the number of nodes grows; CAN scales by convention for how identifiers
are assigned and filtered, and there is no address field to widen.

% ------------------------------------------------------------------------------------------
\section{The CAN frame format}
\label{c10:sec:frameformat}

\lecfig[0.86]{lectures/L10/appendix/images/standard_can_frame_format.png}{CAN's standard data frame,
field by field, in transmission order.}{c10:fig:canframe}

This book implements CAN's \textbf{standard data frame (base format)}, bit by bit as the standard
defines it: an 11-bit identifier and up to 8 data bytes. Remote frames, extended 29-bit identifiers
and everything the standard defines \emph{around} data frames (error and overload frames, and the
3-bit intermission separating one frame from the next) are outside the book's scope, but every frame
this controller puts on the wire is a fully standard-compliant data frame. It is the same
simplification real designs often start from, and it is more than enough to build a genuine, working
controller.

The diagram gives the transmission order and every field's width. Bit stuffing applies to every field
from \textbf{SOF through the CRC field}, and not the CRC delimiter, the ACK field or EOF.

\subsection*{Field by field}
\paragraph{SOF, 1 bit, always dominant}
Sent out on an otherwise idle, recessive bus, so it produces the single recessive-to-dominant edge
every node waits for. That edge marks both that a frame is starting and resynchronises every node's
bit timing against it; \code{bit_timer.resync} in \chapref{c12:ch} is exactly this.

\paragraph{Identifier, 11 bits, MSB first}
Doubles as the message's arbitration priority. Every transmitting node compares every bit it sends
against what the bus actually reads, and drops out the moment it sends recessive but observes
dominant.

\paragraph{RTR, 1 bit, always dominant here}
The Remote Transmission Request bit ends the 12-bit arbitration field: dominant marks a data frame,
recessive a \emph{remote frame}, a request for whoever owns the identifier to send its data. This
book supports data frames only, so its controller always sends RTR dominant, and a receiver sampling
it recessive abandons the frame and sets \code{error} (\chapref{c17:ch} implements that check).
Formally, arbitration spans the identifier \emph{and} RTR; in real CAN that is what lets a data frame
beat a remote frame with the same identifier. With every node here sending RTR dominant, the
identifier always decides on its own.

\paragraph{Control field, 6 bits: IDE, r0, DLC}
\code{IDE} is always \code{0} here; \code{1} would mark an extended-identifier frame, outside the
book's scope. \code{r0} is reserved, always \code{0}. \code{DLC} is the 4-bit-wide length field,
\code{0}--\code{8} bytes encoded as \code{0000}--\code{1000}. It is the length field from
\secref{c10:sec:framing} in a new guise: without it a receiver could not tell whether the data field
ends after this byte or carries on.

\paragraph{Data field, 0 to 8 bytes}
Sent MSB first, byte 0 first, with the length DLC announced. When DLC is \code{0} the field is absent
entirely and the CRC follows directly on the control field.

\paragraph{CRC field, 15 bits plus a delimiter}
The CRC covers every bit from SOF through the end of the data field, and is far more robust than the
summed checksum in \secref{c10:sec:framing}. One detail matters later on: it covers only the
\emph{real} bits. Stuff bits are inserted after the CRC has been computed on the way out and
discarded before it is checked on the way in, so they never enter the computation (\chapref{c14:ch}
and \ref{c15:ch} are where that gating is built). \Chapref{c13:ch}'s engine computes it bit-serially,
one bit at a time while the frame is being sent or received, with no separate ``compute the
checksum'' step at the end.

The \textbf{CRC delimiter} is the 16th bit, always recessive and never stuffed. It exists so that a
bit independent of the bit stuffing rule always follows the CRC value, which is what lets a receiver
find the boundary between the CRC and ACK fields.

\paragraph{ACK field, 2 bits}
The transmitter sends both recessive. Every \emph{other} node that received the frame correctly
(framing intact, stuffing respected, CRC valid) pulls the ACK slot dominant. The transmitter never
checks its own CRC against itself, so the dominant bit from somebody else is its only confirmation
that the frame arrived intact somewhere. The \textbf{ACK delimiter} is recessive and unstuffed, and
marks the boundary for the same reason as the CRC delimiter.

\paragraph{EOF, 7 recessive bits, unstuffed}
Seven identical bits in a row are far more than the bit stuffing rule would ever allow inside a
stuffed field, which is exactly what makes EOF unmistakable. That holds only if every node agrees not
to stuff it, which is why EOF and the two delimiters are defined as exempt from stuffing.

\Canref{2} goes through the same frame field by field, and includes what lies outside this book's
controller: the RTR semantics of remote frames, what the ACK slot means to a transmitter reading it
back, and the three bits of intermission separating one frame from the next.

% ------------------------------------------------------------------------------------------
\section{Bit stuffing}
\label{c10:sec:stuffing}

\textbf{The rule.} Throughout every stuffed field, SOF through the CRC field: after \textbf{five
consecutive bits of the same value} the transmitter inserts an extra bit of the \emph{opposite}
value. That bit carries no data, and every receiver discards it by counting the same way.

\textbf{Why it exists}, for two related reasons:
\begin{enumerate}
  \item \textbf{Clock recovery.} Every node's bit timer runs on its own oscillator. A long run of
    identical bits contains no transitions, so a receiver has nothing confirming that it is still
    aligned with the transmitter's bit boundaries. Guaranteeing a transition at least every fifth
    real bit bounds how far its timing can drift. Real CAN resynchronises on the edges bit stuffing
    guarantees; this book's controller is simpler and resynchronises only at SOF
    (\code{bit_timer.resync}, \chapref{c12:ch}), which makes the bounded drift more important still.
  \item \textbf{Telling real data from EOF.} EOF is seven identical recessive bits, unstuffed. A
    stuffed field can never accidentally produce that run, which is what makes seven unstuffed
    recessive bits an unambiguous end of frame rather than something the identifier or the data field
    could have produced by chance.
\end{enumerate}

\paragraph{Worked example}
To send the 8-bit sequence \code{11111000} the rule fires after the fifth consecutive \code{1}:
insert a \code{0}, then carry on with the real data. Every column below is a bit actually on the
wire, so the real-bits row leaves a gap where the stuff bit sits:

\begin{textdrawing}
Real bits:                1  1  1  1  1  -  0  0  0
Transmitted bit stream:   1  1  1  1  1  0* 0  0  0
Run length:               1  2  3  4  5  1  2  3  4
                                         ^ stuff bit - not real data
\end{textdrawing}

Read the run-length row against the transmitted row, not against the real-bits row: the rule counts
bits on the wire. The stuff bit is itself the first bit of the new run, which is why the count reads
\code{1} under it and \code{2} under the first real \code{0} that follows. \Chapref{c14:ch} builds
exactly this counter in hardware.

A receiver counts the same five \code{1}s, expects the next bit to be a stuff bit, and discards it if
it inverts the run. If a \emph{sixth} consecutive \code{1} turns up instead where the stuff bit
should have been, that is a stuffing violation: a detected framing error, treated as fatal, reported
by \code{rx_shift_reg.stuff_error} in \chapref{c15:ch}.

\begin{warning}
\textbf{A subtlety worth mentioning now, since it surprises people later.} The run count is
\emph{not} cleared at field boundaries. A run beginning in one field and carrying on into the next is
counted straight through, and can trigger a stuff bit right at the boundary. Stuffing cares only
about the bits actually on the wire, never about which named field they belong to.
\end{warning}

\Canref{3} derives the five rule from the clock recovery problem and takes the CRC-15 in the same
chapter, for the reason \chapref{c13:ch} and \ref{c14:ch} then live with: stuffing and the checksum
share a bit stream but not bits, and whoever understands one without the other gets the wrong CRC.

% ------------------------------------------------------------------------------------------
\section{Arbitration, worked through in full}
\label{c10:sec:arbitration}

Nodes P (identifier \code{0x0F0}) and Q (identifier \code{0x100}) start transmitting at the same SOF.
Both send SOF as \code{0}, so no contest arises yet. Then the identifiers, MSB first:

\begin{textdrawing}
             bit:  1  2  3  4  5  6  7  8  9 10 11
 P (0x0F0) sends:  0  0  0  1  1  1  1  0  0  0  0
 Q (0x100) sends:  0  0  1  -  -  -  -  -  -  -  -   (loses at bit 3, stops driving)
       bus reads:  0  0  0  1  1  1  1  0  0  0  0
\end{textdrawing}

Bits 1--2 match for both nodes, so nothing happens. At bit 3, P sends dominant and Q recessive, so
the bus reads dominant. Q, monitoring while it transmits, sees its own \code{1} against a bus reading
\code{0}: \textbf{Q has lost arbitration} and stops driving immediately. P sees \code{0} on a bus
reading \code{0}, exactly what it expected, and carries on without ever finding out that a contest
took place.

In a production controller Q would now receive the rest of P's frame. This book's
\code{can_controller} stops driving at the same moment but abandons the frame, as flagged above and
as \secref{c19:sec:gaps} returns to.

The mechanism generalises to any number of simultaneous transmitters: the numerically lowest
identifier wins (equivalently: the longest run of leading dominant bits), and every loser drops out
the moment its recessive bit is overruled, never later than exactly that bit.

\Canref{4} works through both the two-node and the three-node case, and then takes the two questions
this chapter does not ask: what starvation and priority inversion mean on a bus where low identifiers
always win, and what a transmitter never gets to find out about why it lost.

% ------------------------------------------------------------------------------------------
\section{Bit timing and the sample point}
\label{c10:sec:bittiming}

Arbitration works only if every node agrees on \emph{when} a bit is on the bus, which raises a
question: at what point within a bit period does a node read it?

Every node divides its own clock into bit periods of fixed length using its own local bit timer.
Because those oscillators are independent, and because signals take time to propagate along the bus,
a node does \textbf{not} sample at the start of the period. It waits until a \textbf{sample point}
further into the period, by which time the transmitted level has had time to settle.

This book samples 70~\% into every bit period, a reasonable and simple choice for a short bus on a
single board at 1~Mbit/s.

Real CAN controllers make the sample point configurable and divide the bit period into named segments
(propagation segment, phase segments) tuned against bus length, propagation delay and oscillator
tolerance. That lies outside this book's fixed sample point, but not outside \canref{5}, which
divides the bit period into its segments, works out the timing from a clock frequency and a bit rate,
and describes the resynchronisation at every edge that this controller replaces with a single pulse
at SOF.

\Secref{c10:sec:candef} turns those 70~\% into the concrete VHDL constant \code{SAMPLE_TICK}, derived
from the clock frequency and the bit rate rather than hardcoded.
